\chapter{Einleitung}
\label{ch:einleitung}

\section{Motivation}
Bier ist eines der ältesten Lebensmittel der Menschheit. Spätestens seit dem
Reinheitsgebot von 1516 besteht es in Mitteleuropa aus nur vier Zutaten: Wasser,
Malz, Hopfen und Hefe \citep{hopfner2019reinheit}. Trotz dieser Einfachheit
unterscheiden sich Biere enorm in Geschmack und Geruch -- und ein großer Teil
dieser Vielfalt entsteht nicht im Sudhaus, sondern im Gärtank \citep{malz2021ester}.

Craft-Brauereien experimentieren daher zunehmend mit Gärparametern. Gleichzeitig
fehlt es an \textbf{systematischen Daten}, wie stark einzelne Parameter das
Aroma tatsächlich verändern.

\section{Forschungsfragen}
\begin{enumerate}[label=\textbf{F\arabic*}]
  \item Wie verändert sich die Konzentration aromaaktiver Ester mit der Gärtemperatur?
  \item Ab welcher Temperatur nehmen Verkosterinnen und Verkoster den Unterschied wahr?
  \item Lässt sich das Aromaprofil mit einem einfachen Modell vorhersagen?
\end{enumerate}

\section{Aufbau der Arbeit}
\Cref{ch:grundlagen} erklärt die Biochemie der Gärung, \cref{ch:methodik}
beschreibt die Versuchssude, \cref{ch:ergebnisse} präsentiert Messwerte und
Verkostung, und \cref{ch:fazit} gibt einen Ausblick -- inklusive Rezeptempfehlung.
